\section{Objectifs}
\begin{itemize}
\item Générer des listes.
\item Manipuler éléments et indices.
\item Parcourir une liste et organiser un programme modulaire.
\end{itemize}

\section{Prérequis}
Variables, affectations, conditions, boucles et fonctions Python.

\section{Activité de découverte}
Une suite, un tableau de valeurs ou un échantillon statistique se représente naturellement par une liste.

\section{Vocabulaire}
Liste, élément, indice, longueur, compréhension, parcours.

\section{Cours}
En Première spécialité, la variable, les conditions, les boucles et les fonctions sont consolidées. La seule notion algorithmique nouvelle explicitement annoncée par le BO est la \textbf{liste}.

Une liste peut être écrite en extension, construite par ajouts successifs ou par compréhension :
\begin{verbatim}
L = [k*k for k in range(6)]
\end{verbatim}

\texttt{L[i]} accède à l’élément d’indice \texttt{i} et \texttt{len(L)} donne la longueur. On distingue le parcours des éléments
\begin{verbatim}
for x in L:
    ...
\end{verbatim}
du parcours des indices.

\section{Exemples corrigés}
\texttt{[2*n+1 for n in range(5)]} produit \texttt{[1, 3, 5, 7, 9]}.

\section{Méthode}
Découper une tâche complexe en fonctions courtes : génération, transformation, calcul et restitution.

\section{Erreurs fréquentes}
\begin{itemize}
\item Confondre indice et valeur.
\item Oublier que le premier indice est 0.
\item Introduire d’autres collections comme nouvelles notions : le BO demande de se limiter aux listes.
\end{itemize}

\section{Synthèse}
Les listes donnent une structure commune aux suites, tableaux de valeurs, séries statistiques et simulations.
